\labsection{7}{Supply Chain Management: Procure, Receive, and Measure}{sec:lab07}

\begin{labproject}{Build the upstream supply chain for Road Bikes and connect it to SCM performance}
\textbf{Coverage.} Tutorial 7 + Chapter 7.\\
\textbf{Core system.} Odoo Purchase + Inventory + Manufacturing.\\
\textbf{Goal.} Demonstrate supplier master data, purchase orders, receipts, inventory availability, and the link between supplier performance and production readiness.

\begin{labmode}
Tutorial 7 asks students to research an SCM implementation and compare CRM with SCM. The Odoo work gives a concrete supply-chain process you can analyse using metrics such as lead time, fill rate, and on-time performance.
\end{labmode}

\medskip\noindent\textbf{Part A --- Create vendors and vendor-product links.}

Create the supplied vendors and link them to raw materials on the product Purchase tab:
\begin{itemize}
  \item Ipoh Metal Supply $\rightarrow$ Aluminium Frame;
  \item RubberWheel Co. $\rightarrow$ Rubber Tires;
  \item BikeParts Supplier $\rightarrow$ Brake Set, Gear Set, Bicycle Chain;
  \item CarbonTech Components $\rightarrow$ Carbon Fiber Frame if used in your extended catalogue.
\end{itemize}

Set vendor prices from the guide and record a delivery lead time if the field is present.

\medskip\noindent\textbf{Part B --- Raise RFQs and confirm Purchase Orders.}

Use one RFQ per vendor. Assign the Purchasing Officer as Buyer. Use the guide's example quantities or quantities approved by the tutor. Confirm each RFQ so the document becomes a Purchase Order and a Receipt smart button appears.

\medskip\noindent\textbf{Part C --- Receive materials.}

Open the linked Receipt, check quantities, and Validate. Repeat for each supplier. Verify on-hand stock for one component before and after validation.

\medskip\noindent\textbf{Part D --- Explain why receipt is not the same as order.}

Record one sentence for each:
\begin{itemize}
  \item what the Purchase Order proves;
  \item what the Receipt proves;
  \item which event makes components usable by Manufacturing;
  \item which event should supplier on-time/fill-rate analysis compare against.
\end{itemize}

\medskip\noindent\textbf{Part E --- Build a supplier scorecard concept.}

For each supplier, define:
\begin{itemize}
  \item initial fill rate;
  \item order lead time;
  \item on-time performance;
  \item one quality-related measure of your own.
\end{itemize}

You do not need a large history in the demo database. The purpose is to define how repeated Odoo Purchase/Receipt transactions could generate the metric over time.

\medskip\noindent\textbf{Part F --- Complete Tutorial 7 research.}

Group 1: research a manufacturing company that successfully implemented SCM, whether ERP existed before SCM, the challenges, and how success was measured.\\
Group 2: compare CRM and SCM, including similarities, differences, interaction points, and technologies.

Use the SCM metrics from Chapter 7 to strengthen the Group 1 analysis and the demand-to-fulfilment hand-off to strengthen the Group 2 comparison.
\end{labproject}

\begin{labdeliverable}
Submit/share the Tutorial 7 presentation and reflection as instructed. Include Odoo evidence of one vendor-product link, one confirmed Purchase Order, and one Done Receipt. Your reflection should connect at least one research point back to the Odoo supply-chain process.
\end{labdeliverable}
